\subsection{Group of types of quadratic forms.}

We shall equip the set \(\mathfrak{W}\) of types of nondegenerate quadratic forms of finite dimensions over A with a structure of a commutative group. We will define an addition in \(\mathfrak{W}\) by the formula

\[
\mathrm{T} + \mathrm{T} ^ {\prime} = \theta (\mathrm{T} _ {\tau} \mathrm{T} ^ {\prime})\tag{2}
\]

\[
(\mathrm{T}, \mathrm{T} ^ {\prime} \text {   in   } \mathfrak {W}),
\]

This addition is commutative since \(T^{\prime} \tau T\) is equivalent to \(T \tau T^{\prime}\) . It is associative because, if T, \(T^{\prime}\) and \(T^{\prime\prime}\) are elements of W, we have

\[
\begin{array}{r l} & (\mathrm{T} + \mathrm{T} ^ {\prime}) + \mathrm{T} ^ {\prime \prime} \sim (\mathrm{T} + \mathrm{T} ^ {\prime}) \tau \mathrm{T} ^ {\prime \prime} \sim (\mathrm{T} \tau \mathrm{T} ^ {\prime}) \tau \mathrm{T} ^ {\prime \prime} \\ & \qquad \sim \mathrm{T} \tau (\mathrm{T} ^ {\prime} \tau \mathrm{T} ^ {\prime \prime}) \sim \mathrm{T} \tau (\mathrm{T} ^ {\prime} + \mathrm{T} ^ {\prime \prime}) \sim \mathrm{T} + (\mathrm{T} ^ {\prime} + \mathrm{T} ^ {\prime \prime}), \end{array}
\]

whence \((\mathrm{T} + \mathrm{T}' ) + \mathrm{T}'' = \mathrm{T} + (\mathrm{T}' + \mathrm{T}'' )\) since two elements of W that have the same type are equal. Moreover, the addition just defined has an identity element: indeed, it is clear that all neutral forms have the same type \(T_{0}\) , namely that of the zero form of dimension zero; one sees immediately that \(T_{0}\) is the sought neutral element. Finally, the existence, for every \(T \in W\) , of an element opposite to T follows immediately from the following proposition:

PROPOSITION 3. — Let Q be a nondegenerate quadratic form of finite dimension on a vector space V over A. Denote by −Q the quadratic form on V defined by \((-Q)(x) = -Q(x) (x \in V)\) . Then the form \(Q \tau (-Q)\) is neutral.

Indeed, the restriction of \(Q \tau (-Q)\) to the diagonal \(D\) of \(V \times V\) is zero. The index of this form is therefore \(\geqslant\frac{1}{2}\dim(V\times V)\) ( \(\S\) 4, no 2, def. 2), and consequently is equal to \(\frac{1}{2}\dim(V\times V)\) (ibid., formula (4)). It follows that \(Q\tau(-Q)\) is neutral (ibid.).

This allows us to state the following definition:

DEFINITION 1. --- The set of types of nondegenerate quadratic forms of finite dimension over A, equipped with the addition defined by (2), is called the group of types of quadratic forms, or Witt group, of A.

Remarks. --- 1) Every nondegenerate quadratic form Q of finite dimension whose type is zero (that is, such that \(\theta(Q) = T_0\) with the notations above) is a neutral form. Indeed, there exist neutral forms N, N′ such that \(Q \tau N\) is equivalent to N′. This shows that Q is of even dimension, hence that there exists a neutral form \(N_1\) of the same dimension as Q. Since Q and \(N_1\) have the same type, it follows from prop. 1 that they are equivalent, hence that Q is neutral.

\begin{enumerate}
\def\labelenumi{\arabic{enumi})}
\setcounter{enumi}{1}
\tightlist
\item
  For every quadratic form Q of finite dimension over A, denote by \(\delta(Q)\) the class modulo 2 of the dimension of Q. We have
\end{enumerate}

\[
\delta (Q \tau Q ^ {\prime}) = \delta (Q) + \delta (Q ^ {\prime}).
\]

Since every neutral form N has even dimension, we have \(\delta(N) = 0\) ; the relation \(Q \sim Q'\) therefore implies \(\delta(Q) = \delta(Q')\) . Thus the restriction of \(\delta\) to the group \(\mathfrak{W}\) of types of quadratic forms on A is a homomorphism from \(\mathfrak{W}\) into the group \(\mathbf{Z}/(2)\) . This homomorphism is surjective when A has characteristic \(\neq 2\) , but is not so if A has characteristic 2, because a quadratic form of odd dimension is then degenerate since its associated bilinear form is alternating (cf.~§ 5).

\begin{enumerate}
\def\labelenumi{\arabic{enumi})}
\setcounter{enumi}{2}
\tightlist
\item
  Let \(a\) be an element \(\neq 0\) of A. If N is a neutral form, then so is \(aN\) . It follows that the relation \(Q \sim Q'\) implies \(aQ \sim aQ'\) . For every element T of the group \(\mathfrak{W}\) , we set
\end{enumerate}

\[
a. \mathrm{T} = \theta (a \mathrm{T}).\tag{3}
\]

We thus obtain an external law of composition between the group A* of nonzero elements of A and the group W. The following formulas follow immediately from the definition:

\[
a. (\mathrm{T} + \mathrm{T} ^ {\prime}) = a. \mathrm{T} + a. \mathrm{T} ^ {\prime}, \quad a b. \mathrm{T} = a. (b. \mathrm{T})\tag{4}
\]

\[
(a, b \text {   in   } A ^ {*}, T, T ^ {\prime} \text {   in   } \mathfrak {W}).
\]

On the other hand, if \(a\) , \(b\) and \(a + b\) are in \(\mathbf{A}^*\) one does not in general have \((a + b)\) . \(\mathrm{T} = a\) . \(\mathrm{T} + b\) . \(\mathrm{T}\) ( \(\mathrm{T} \in \mathfrak{W}\) ).

PROPOSITION 4. --- Let Q be a nondegenerate quadratic form on a finite-dimensional vector space E over A. Suppose A has characteristic \(\neq 2\) , and let \((x_{1},\ldots,x_{n})\) be an orthogonal basis of V. Let \(T_{1}\) denote the type of the quadratic form \(Q_{1}\) defined on the vector space A and such that \(Q_{1}(1)=1\) . The type of Q is then \(\sum_{i=1}^{n}\mathrm{Q}(x_{i}).T_{1}\) .

Indeed, the form Q is equivalent to

\[
(\mathrm{Q} (x _ {1}) \mathrm{Q} _ {1}) \tau \dots \tau (\mathrm{Q} (x _ {n}) \mathrm{Q} _ {1})
\]

COROLLARY. --- With the hypotheses and notations being those of prop. 3, the elements \(a\) . \(\mathrm{T}_{1}\) ( \(a \in \mathrm{A}^{*}\) ) form a set of generators of the group of types of quadratic forms on A.

Thus, determining the structure of the group of types of quadratic forms on A amounts to finding the Z-linear relations that exist among the elements of the form \(a.T_{1}\) . If \(b \in A^{*}\) , the form \(Q_{1}\) defined in Prop. 4 is clearly equivalent to \(b^{2}Q_{1}\) ; one therefore has \(a.T_{1} = ab^{2}.T_{1}\) , which shows that \(a.T_{1}\) depends only on the class of a modulo the subgroup \((A^{*})^{2}\) of squares of elements of \(A^{*}\) . Moreover, it follows from prop. 3 that one has \((-a).T_{1} = -a.T_{1}\) . However, there exist in general other Z-linear relations among the \(a.T_{1}\) besides those which are deduced from the relations we have just indicated.

PROPOSITION 5. --- Suppose that A is a maximal ordered field. Let Q be a nondegenerate quadratic form of finite dimension over A, and let (s, t) be its signature (§ 7, no. 2, def. 2). Then the type of Q is (s - t)·T₁, and the group W of types of quadratic forms over A is an infinite cyclic group generated by T₁.

Indeed, as \(\mathbf{A}^{*}/(\mathbf{A}^{*})^{2}\) has order 2 and that \((-1)\) . \(\mathrm{T}_{1} = -\mathrm{T}_{1}\) , \(\mathfrak{W}\) is generated by \(\mathrm{T}_{1}\) and is consequently monogenic. For all \(n > 0\) , \(n\) . \(\mathrm{T}_{1}\) is the type of positive nondegenerate quadratic forms of dimension \(n\) ; since these forms are not neutral, we have \(n\) . \(\mathrm{T}_{1} \neq 0\) , which shows that \(\mathfrak{W}\) is infinite. Finally, a form of signature \((s, t)\) is isomorphic, with the notations of Prop. 4, to the direct sum of s forms \(Q_{1}\) and of t forms \(-Q_{1}\) (§ 7, no. 2, th. 1); it follows that its type is \((s - t)\) . \(T_{1}\) .

